\chapter{问题定义和代表性方法}
\label{chap:problem}

\section{POI推荐问题定义}
\label{sec:poi_problem}

\subsection{用户-POI交互序列}
\label{sec:user_poi_sequence}

给定用户 $u$ 的历史签到序列：

\begin{equation}
S_u = [(poi_1, t_1), (poi_2, t_2), \dots, (poi_n, t_n)]
\label{eq:user_sequence}
\end{equation}

其中 $S_u$ 表示用户 $u$ 的历史交互序列，$poi_i$ 为第 $i$ 次交互的POI，$t_i$ 为对应的时间戳。

\subsection{Next-POI预测任务}
\label{sec:next_poi_task}

基于用户历史序列和上下文信息，预测用户下一个可能访问的POI：

\begin{equation}
poi_{n+1} = \arg\max_{poi \in \mathcal{P}} P(poi \mid S_u, context)
\label{eq:next_poi}
\end{equation}

其中 $\mathcal{P}$ 为所有POI的集合，$context$ 包含时空上下文信息。

\section{语义ID的定义与表示}
\label{sec:sid_definition}

本文提出将POI映射为离散的语义标识符（Semantic ID），使其在保留语义相关性的同时具备可枚举的离散表示。

\subsection{层次化量化结构}
\label{sec:hierarchical_quant}

语义ID采用多层量化编码结构，从粗到细逐层区分POI：

\begin{equation}
SID(poi) = \langle a_{p1} \rangle \langle b_{p2} \rangle \langle c_{p3} \rangle
\label{eq:sid_format}
\end{equation}

例如：$\langle a_{12} \rangle \langle b_{34} \rangle \langle c_{56} \rangle$

其中 $\langle a_X \rangle$、$\langle b_X \rangle$、$\langle c_X \rangle$ 分别对应P1、P2、P3三层量化索引。三层量化使得具有相似语义特征的POI在ID空间中彼此靠近，从而支持基于ID序列的推荐学习。

各层的设计意图如下：P1层捕获宏观特征（如地理大区、核心业态）；P2层细化到城市级别特征；P3层编码细粒度差异（如同一区域内不同类别）。三层码本相互独立，协同实现语义粒度的层次化表征。

\subsection{冲突处理原则}
\label{sec:collision_principle}

由于离散码本容量有限，多个POI可能被量化到相同的语义ID。此时需要附加额外信息进行消歧，确保每个POI拥有全局唯一的标识符。消歧策略将在第三章中详述。

\section{代表性方法概述}
\label{sec:related_methods}

本章概述与本文研究密切相关的几类代表性方法，包括序列推荐、语义编码与大语言模型推荐三个方向。

\subsection{序列推荐方法}
\label{sec:sequential_rec}

序列推荐旨在利用用户历史行为序列建模其动态偏好。早期代表性方法包括马尔可夫链模型（如FPMC~\cite{rendle2010factorifying}），通过对用户-物品交互矩阵进行分解并引入马尔可夫链建模序列转移。后续研究在此基础上引入注意力机制（如STAN~\cite{wang2021spatial}）以建模更复杂的序列依赖关系。

深度学习方法进一步推动了序列推荐的发展。RNN类方法（如GRU4Rec~\cite{hidasi2015session}）利用循环神经网络建模长序列依赖；Transformer类方法（如SASRec~\cite{kang2018self}）通过自注意力机制实现更高效的位置建模；时空注意力网络（如STAN）则在注意力基础上显式融合空间和时间信息。

上述方法均以预测下一个交互物品为直接目标，在隐式表示空间中运作，难以直接解释推荐理由。

\subsection{语义编码方法}
\label{sec:semantic_encoding}

语义编码方法的核心思想是将物品映射到低维语义向量或离散代码，以支持可解释的推荐推理。

VGAE~\cite{kipf2016variational}和VAE系列方法利用变分自编码器学习物品的连续语义表示；向量量化技术（如VQ-VAE~\cite{van2017neural}）进一步将连续表示离散化为紧凑码字，支持基于码本的物品检索。

近年来，研究者开始关注语义标识符（Semantic ID）的设计，旨在生成兼具语义相关性和离散可枚举性的物品代码。在推荐场景中，语义ID有望替代传统物品ID，支持基于ID序列的生成式推荐~\cite{zhou2024sid}。

本文借鉴上述思路，提出基于残差量化变分自编码器（RQ-VAE）的语义ID生成方法，并将其作为大语言模型推荐的输入序列。

\subsection{大语言模型推荐}
\label{sec:llm_rec}

大语言模型（LLM）凭借其强大的语义理解和推理能力，在推荐系统中展现了新的潜力。近期研究表明，LLM可在零样本或微调条件下完成推荐任务，包括next-item预测~\cite{heidari2024llm4rec}、序列推荐~\cite{liu2024llmseq}等。

指令微调（Instruction Tuning）是实现LLM推荐的主流范式：通过设计包含用户信息、历史行为和推荐目标的Prompt模板，将推荐任务转化为条件文本生成任务。代表性工作包括InstructRec~\cite{zhange2023instructrec}和LLM-Rec~\cite{liu2024llm4rec}等。

然而，将LLM直接应用于推荐仍面临挑战：传统物品ID缺乏语义信息，难以被LLM理解；直接使用物品属性文本作为标识则序列过长且检索效率低。为解决上述问题，本文提出基于语义ID的Prompt设计策略，并通过QLoRA微调使LLM适配POI推荐任务。

\section{数据划分策略}
\label{sec:data_split}

\subsection{Yelp数据集城市分布统计}
\label{sec:city_distribution}

Yelp数据集覆盖24个州、1,261个城市，共117,788个POI。城市POI数量分布呈长尾分布，如表所示：

\begin{table}[htbp]
\centering
\caption{Yelp数据集Top-10城市POI数量统计}
\label{tab:city_distribution}
\begin{tabular}{llr}
\toprule
排名 & 城市 & POI数量 \\
\midrule
1 & Philadelphia (PA) & 11,711 \\
2 & Tampa (FL) & 6,943 \\
3 & Tucson (AZ) & 6,701 \\
4 & Indianapolis (IN) & 6,010 \\
5 & Nashville (TN) & 5,554 \\
6 & New Orleans (LA) & 5,155 \\
7 & Edmonton (AB) & 4,381 \\
8 & Reno (NV) & 4,143 \\
9 & Saint Louis (MO) & 3,795 \\
10 & Santa Barbara (CA) & 2,826 \\
\bottomrule
\end{tabular}
\end{table}

\subsection{Philadelphia数据子集构建规则}
\label{sec:philly_data_rules}

当针对Philadelphia城市构建训练数据子集时，采用以下规则：

\textbf{1. 城市级数据抽取}：从按州/城市分层的目录结构（`PA/Philadelphia/`）中直接加载该城市的全部POI数据，避免跨城市数据污染。

\textbf{2. 特征维度设计}：采用GNPR-SID V2特征方案，总维度79维（无属性）或143维（含属性），具体构成为：
\begin{itemize}
    \item 类别嵌入：64维，使用SentenceTransformer或可学习哈希嵌入
    \item 空间坐标：3维，经纬度转换为单位球面3D坐标
    \item 时间特征：12维，Fourier展开捕捉周期性（6频率$\times$sin/cos）
    \item 属性嵌入（可选）：64维，从POI的attributes字段文本编码
\end{itemize}

\textbf{3. 数据质量过滤}：仅保留`is_poi=True`的记录，确保为有效商户而非其他类型实体。

\textbf{4. 时间分布聚合}：从用户签到记录（review_poi.json）中统计每个POI的访问时间分布，包括时段比例（早/午/晚/夜）和周内比例（工作日/周末）。

\textbf{5. 类别标签处理}：从POI的categories字段提取所有类别标签，映射到类别词表。单个POI最多保留5个类别，不足部分补零。

Philadelphia数据子集共11,711个POI，675,543条用户评论，为训练语义ID模型提供充足的样本量。

\subsection{选择Philadelphia作为训练数据的原因}
\label{sec:philly_selection}

\textbf{选择Philadelphia作为训练数据的原因}：
\begin{enumerate}
    \item \textbf{数据规模最大}：11,711个POI，675,543条用户评论（占全量数据集的14.9\%），数据量充足
    \item \textbf{避免跨城市复杂性}：单一城市区域，POI之间的地理距离和语义相关性更加一致
\end{enumerate}

